\documentclass[aspectratio=169,10pt]{beamer}

% ------------------------------------------------------------
% Curso Completo: Aprendizaje por Refuerzo
% "De la Probabilidad a las Redes Neuronales Profundas"
% Formato: Diapositivas de Alta Definición (16:9)
% Compilar con pdfLaTeX
% ------------------------------------------------------------

\usepackage[T1]{fontenc}
\usepackage[utf8]{inputenc}
\usepackage[spanish,es-nodecimaldot]{babel}
\usepackage{lmodern}
\usepackage{microtype}
\usepackage{amsmath,amssymb,mathtools}
\usepackage{booktabs}
\usepackage{listings}
\usepackage{tikz}
\usetikzlibrary{arrows.meta,calc,positioning,shapes.geometric}

% ---------- Paleta de alto contraste ----------
\definecolor{Ink}{HTML}{102A43}       % azul tinta: texto principal
\definecolor{Paper}{HTML}{FFFFFF}     % fondo obligatorio
\definecolor{Mist}{HTML}{D9E2EC}      % gris azulado para guías y divisores
\definecolor{Teal}{HTML}{005EA8}      % azul accesible: color primario
\definecolor{TealDark}{HTML}{003E73}  % azul profundo
\definecolor{Coral}{HTML}{D9480F}     % naranja oscuro: énfasis
\definecolor{Gold}{HTML}{B77900}      % ámbar oscuro: recompensa
\definecolor{Leaf}{HTML}{2F6F44}      % verde oscuro: resultados positivos
\definecolor{CodeBg}{HTML}{102A43}    % código con contraste alto
\definecolor{SoftText}{HTML}{486581}  % texto secundario legible

% ---------- Identidad visual Beamer ----------
\usefonttheme{professionalfonts}
\setbeamercolor{normal text}{fg=Ink,bg=Paper}
\setbeamercolor{structure}{fg=Teal}
\setbeamercolor{frametitle}{fg=Ink,bg=Paper}
\setbeamercolor{alerted text}{fg=Coral}
\setbeamercolor{itemize item}{fg=Coral}
\setbeamercolor{itemize subitem}{fg=Teal}
\setbeamercolor{block title}{fg=white,bg=Teal}
\setbeamercolor{block body}{fg=Ink,bg=Teal!8}
\setbeamercolor{block title alerted}{fg=white,bg=Coral}
\setbeamercolor{block body alerted}{fg=Ink,bg=Coral!8}
\setbeamercolor{block title example}{fg=white,bg=Leaf}
\setbeamercolor{block body example}{fg=Ink,bg=Leaf!9}
\setbeamercolor{background canvas}{bg=Paper}

\setbeamertemplate{navigation symbols}{}
\setbeamertemplate{blocks}[rounded][shadow=false]
\setbeamertemplate{itemize item}{\raise0.15ex\hbox{\tikz\fill[Coral] (0,0) circle (0.55ex);}}
\setbeamertemplate{itemize subitem}{\raise0.15ex\hbox{\tikz\fill[Teal] (0,0) circle (0.4ex);}}
\setlength{\leftmargini}{1.5em}

% Barra superior de título
\setbeamertemplate{frametitle}{%
  \nointerlineskip
  \begin{beamercolorbox}[wd=\paperwidth,ht=0.95cm,dp=0.25cm,leftskip=0.55cm,rightskip=0.4cm]{frametitle}
    \usebeamerfont{frametitle}\insertframetitle
    \hfill
    \tikz[baseline=-0.5ex]{
      \fill[Teal] (0,0) circle (0.10);
      \fill[Coral] (0.28,0) circle (0.10);
      \fill[Gold] (0.56,0) circle (0.10);
    }
  \end{beamercolorbox}
  \vspace{-0.12cm}
  \hspace{0.55cm}\color{Teal}\rule{1.25cm}{1.4pt}%
}

% Pie con sección y barra de progreso
\newcommand{\progressbar}{%
  \begin{tikzpicture}[baseline]
    \fill[Mist] (0,0) rectangle (3.0,0.055);
    \pgfmathsetmacro{\progress}{3.0*\insertframenumber/max(1,\inserttotalframenumber)}
    \fill[Teal] (0,0) rectangle (\progress,0.055);
  \end{tikzpicture}%
}
\setbeamertemplate{footline}{%
  \leavevmode
  \begin{beamercolorbox}[wd=\paperwidth,ht=0.42cm,dp=0.14cm,leftskip=0.55cm,rightskip=0.55cm]{normal text}
    \scriptsize\color{SoftText}\insertsectionhead
    \hfill\progressbar\hspace{0.6em}\insertframenumber\,/\,\inserttotalframenumber
  \end{beamercolorbox}
}

% ---------- Comandos pedagógicos del curso ----------
\newcommand{\state}[1]{\tikz[baseline=(x.base)]\node[rounded corners=2pt,fill=Teal!13,text=TealDark,inner xsep=5pt,inner ysep=2pt] (x) {\scriptsize\bfseries ESTADO\enspace #1};}
\newcommand{\rlaction}[1]{\tikz[baseline=(x.base)]\node[rounded corners=2pt,fill=Coral!13,text=Coral,inner xsep=5pt,inner ysep=2pt] (x) {\scriptsize\bfseries ACCIÓN\enspace #1};}
\newcommand{\reward}[1]{\tikz[baseline=(x.base)]\node[rounded corners=2pt,fill=Gold!22,text=Ink,inner xsep=5pt,inner ysep=2pt] (x) {\scriptsize\bfseries RECOMPENSA\enspace #1};}
\newcommand{\keyidea}[1]{%
  \begin{beamercolorbox}[rounded=true,sep=0.60em,wd=\linewidth]{block body alerted}
    \textcolor{Coral}{\bfseries Idea clave}\quad #1
  \end{beamercolorbox}%
}
\newcommand{\module}[2]{%
  \begin{frame}[plain]
    \begin{tikzpicture}[remember picture,overlay]
      \fill[Paper] (current page.south west) rectangle (current page.north east);
      \fill[Teal] (current page.south west) rectangle ([xshift=0.085\paperwidth]current page.north west);
      \fill[Coral] ([xshift=0.085\paperwidth,yshift=0.8cm]current page.south west)
        rectangle ([xshift=0.097\paperwidth,yshift=-0.8cm]current page.north west);
      \draw[Mist,line width=1pt] ([xshift=0.15\paperwidth]current page.west)
        -- ([xshift=-0.8cm]current page.east);
      \node[anchor=west,text=Ink,font=\fontsize{24}{28}\selectfont\bfseries,
            text width=0.72\paperwidth] at ([xshift=0.17\paperwidth,yshift=0.45cm]current page.west) {#1};
      \node[anchor=west,text=SoftText,font=\large,text width=0.68\paperwidth]
        at ([xshift=0.17\paperwidth,yshift=-1.0cm]current page.west) {#2};
      \node[rotate=90,text=white,font=\bfseries\small] at
        ([xshift=0.0425\paperwidth]current page.west) {MÓDULO};
      \fill[Teal] ([xshift=-1.55cm,yshift=0.38cm]current page.east) circle (0.10cm);
      \fill[Coral] ([xshift=-1.22cm,yshift=0.38cm]current page.east) circle (0.10cm);
      \fill[Gold] ([xshift=-0.89cm,yshift=0.38cm]current page.east) circle (0.10cm);
    \end{tikzpicture}
  \end{frame}%
}

% ---------- Configuración de código Python ----------
\lstdefinestyle{rlcode}{
  language=Python,
  backgroundcolor=\color{CodeBg},
  basicstyle=\ttfamily\scriptsize\color{white!92},
  keywordstyle=\bfseries\color{Gold},
  commentstyle=\itshape\color{Leaf!85},
  stringstyle=\color{Coral!85},
  numbers=left,
  numberstyle=\tiny\color{white!35},
  numbersep=7pt,
  frame=none,
  showstringspaces=false,
  breaklines=true,
  tabsize=2,
  xleftmargin=1.3em,
  aboveskip=0pt,
  belowskip=0pt
}

% ---------- Metadatos ----------
\title{Aprendizaje por Refuerzo}
\subtitle{De la Probabilidad a las Redes Neuronales Profundas (36 Horas)}
\author{Curso de Especialización en RL}
\institute{Metodología: Lápiz, Papel y Código · Caso Conductor: OptiBot}
\date{2026}

\begin{document}

% ============================================================
% PORTADA
% ============================================================
\begin{frame}[plain]
  \begin{tikzpicture}[remember picture,overlay]
    \fill[Paper] (current page.south west) rectangle (current page.north east);
    \draw[Ink,line width=1.4pt,rounded corners=7pt]
      ([xshift=0.72\paperwidth,yshift=0.7cm]current page.south west)
      rectangle ([xshift=-0.7cm,yshift=-0.7cm]current page.north east);
    \fill[Teal] ([xshift=0.68\paperwidth]current page.south west)
      -- ([xshift=0.74\paperwidth]current page.south west)
      -- ([xshift=0.68\paperwidth,yshift=0.25\paperheight]current page.south west) -- cycle;

    % Trayectoria abstracta de un agente
    \draw[SoftText!55,line width=1.2pt,dashed]
      ([xshift=-3.6cm,yshift=-2.3cm]current page.north east)
      .. controls +(0.2,-1.1) and +(-1.0,0.5) ..
      ([xshift=-1.2cm,yshift=-4.2cm]current page.north east)
      .. controls +(0.4,-0.8) and +(0.5,0.8) ..
      ([xshift=-2.4cm,yshift=-6.4cm]current page.north east);
    \foreach \x/\y/\c in {-3.6/-2.3/Teal,-1.2/-4.2/Coral,-2.4/-6.4/Gold}{
      \fill[\c] ([xshift=\x cm,yshift=\y cm]current page.north east) circle (0.13cm);
      \draw[Ink,line width=1pt] ([xshift=\x cm,yshift=\y cm]current page.north east) circle (0.23cm);
    }

    \node[anchor=west,text=Teal,font=\bfseries\small] at ([xshift=0.75cm,yshift=2.55cm]current page.west)
      {CURSO COMPLETO TEORÍA + PRÁCTICA (36 HORAS)};
    \node[anchor=west,text=Ink,font=\fontsize{26}{30}\selectfont\bfseries,text width=0.60\paperwidth]
      at ([xshift=0.75cm,yshift=1.15cm]current page.west) {\inserttitle};
    \node[anchor=west,text=SoftText,font=\large,text width=0.55\paperwidth]
      at ([xshift=0.75cm,yshift=-0.15cm]current page.west) {\insertsubtitle};
    \node[anchor=west,text=Ink,font=\small]
      at ([xshift=0.75cm,yshift=-2.35cm]current page.west) {\textbf{\insertauthor}\quad\textcolor{SoftText}{\insertinstitute}};
  \end{tikzpicture}
\end{frame}

% ============================================================
% MAPA DEL CURSO
% ============================================================
\begin{frame}{Ruta de Aprendizaje del Curso}
  \centering
  \resizebox{0.95\textwidth}{!}{%
  \begin{tikzpicture}[node distance=0.35cm and 0.25cm,
    card/.style={rounded corners=4pt,minimum width=2.45cm,minimum height=2.2cm,
      text width=2.15cm,align=left,inner sep=6pt},
    arrow/.style={-{Latex[length=2mm]},line width=1pt,draw=SoftText!60}]
    \node[card,fill=Teal!12] (m0) {\textcolor{Teal}{\bfseries M0 (3h)}\\[0.2em]\bfseries Probabilidad\\[-0.1em]{\scriptsize Incertidumbre y esperanza}};
    \node[card,fill=Coral!11,right=of m0] (m1) {\textcolor{Coral}{\bfseries M1 (3h)}\\[0.2em]\bfseries Bandidos\\[-0.1em]{\scriptsize Explorar vs explotar}};
    \node[card,fill=Gold!18,right=of m1] (m2) {\textcolor{Gold!70!Ink}{\bfseries M2 (3h)}\\[0.2em]\bfseries Cadenas Markov\\[-0.1em]{\scriptsize Dinámica pasiva}};
    \node[card,fill=Leaf!15,right=of m2] (m3) {\textcolor{Leaf!75!Ink}{\bfseries M3 (5h)}\\[0.2em]\bfseries MDP + Bellman\\[-0.1em]{\scriptsize Control y valor}};
    \node[card,fill=Teal!12,right=of m3] (m4) {\textcolor{Teal}{\bfseries M4 (4h)}\\[0.2em]\bfseries Prog. Dinámica\\[-0.1em]{\scriptsize Solución con modelo}};
    \draw[arrow] (m0) -- (m1); \draw[arrow] (m1) -- (m2); \draw[arrow] (m2) -- (m3); \draw[arrow] (m3) -- (m4);
  \end{tikzpicture}%
  }

  \vspace{0.35cm}
  \resizebox{0.80\textwidth}{!}{%
  \begin{tikzpicture}[node distance=0.35cm and 0.25cm,
    card/.style={rounded corners=4pt,minimum width=2.45cm,minimum height=2.2cm,
      text width=2.15cm,align=left,inner sep=6pt},
    arrow/.style={-{Latex[length=2mm]},line width=1pt,draw=SoftText!60}]
    \node[card,fill=Coral!11] (m5) {\textcolor{Coral}{\bfseries M5 (4h)}\\[0.2em]\bfseries Monte Carlo\\[-0.1em]{\scriptsize Episodios completos}};
    \node[card,fill=Gold!18,right=of m5] (m6) {\textcolor{Gold!70!Ink}{\bfseries M6 (6h)}\\[0.2em]\bfseries TD Learning\\[-0.1em]{\scriptsize SARSA y Q-Learning}};
    \node[card,fill=Leaf!15,right=of m6] (m7) {\textcolor{Leaf!75!Ink}{\bfseries M7 (5h)}\\[0.2em]\bfseries Deep RL (DQN)\\[-0.1em]{\scriptsize Redes neuronales}};
    \node[card,fill=TealDark!18,right=of m7] (m8) {\textcolor{TealDark}{\bfseries M8 (3h)}\\[0.2em]\bfseries Proyecto Final\\[-0.1em]{\scriptsize Almacén adaptativo}};
    \draw[arrow] (m5) -- (m6); \draw[arrow] (m6) -- (m7); \draw[arrow] (m7) -- (m8);
  \end{tikzpicture}%
  }

  \vspace{0.35cm}
  \keyidea{Un único hilo conductor: \textbf{OptiBot en el Almacén Inteligente}, resolviendo problemas cada vez más desafiantes.}
\end{frame}

% Inclusión de módulos auxiliares
% ============================================================
% MÓDULO 0: PROBABILIDAD Y PENSAMIENTO ESTOCÁSTICO PARA RL
% ============================================================
\section{M0 · Probabilidad}
\module{Módulo 0: Probabilidad}{El lenguaje de la incertidumbre: de las frecuencias a la esperanza matemática}

\begin{frame}{¿Por qué probabilidad antes de Reinforcement Learning?}
  \begin{columns}[T,onlytextwidth]
    \column{0.48\textwidth}
      \begin{block}{El problema en el mundo real}
        \begin{itemize}
          \item Los robots se resbalan en el suelo húmedo.
          \item Los sensores devuelven mediciones con ruido.
          \item Los tiempos de entrega nunca son exactamente iguales.
        \end{itemize}
      \end{block}
      \vspace{0.2cm}
      \small En el almacén de \textbf{OptiBot}, una misma acción no siempre produce el mismo resultado.
    \column{0.48\textwidth}
      \begin{alertblock}{La solución: Pensamiento Estocástico}
        \begin{itemize}
          \item No podemos predecir el futuro exacto.
          \item \textbf{Sí} podemos modelar la distribución de futuros posibles.
          \item Toda función de valor en RL es una \textbf{esperanza matemática}.
        \end{itemize}
      \end{alertblock}
  \end{columns}
  \vspace{0.4cm}
  \keyidea{En RL no buscamos certezas absolutas, sino maximizar el \textbf{valor esperado} de nuestras ganancias a largo plazo.}
\end{frame}

\begin{frame}{Variables Aleatorias y Distribuciones}
  \begin{columns}[c,onlytextwidth]
    \column{0.50\textwidth}
      \begin{block}{Definición intuitiva}
        Una \textbf{variable aleatoria discreta} $X$ es una regla que asigna un número a cada posible resultado de un experimento.
      \end{block}
      \vspace{0.25cm}
      \textbf{Propiedades obligatorias:}
      \begin{itemize}
        \item $0 \le P(X = x) \le 1$ para todo valor $x$.
        \item $\sum_x P(X = x) = 1$ (la certeza total).
      \end{itemize}
    \column{0.46\textwidth}
      \centering
      \textbf{Caso OptiBot: Tiempo de Entrega $T$}\\
      \vspace{0.2cm}
      \begin{tabular}{@{}ccc@{}}
        \toprule
        \textbf{Días ($t$)} & \textbf{$P(T = t)$} & \textbf{Histograma} \\
        \midrule
        1 día & $0.30$ & \tikz{\fill[Teal] (0,0) rectangle (1.2,0.25);} \\
        2 días & $0.50$ & \tikz{\fill[TealDark] (0,0) rectangle (2.0,0.25);} \\
        3 días & $0.20$ & \tikz{\fill[Coral] (0,0) rectangle (0.8,0.25);} \\
        \bottomrule
      \end{tabular}\\
      \vspace{0.2cm}
      \small Suma: $0.30 + 0.50 + 0.20 = 1.00$
  \end{columns}
\end{frame}

\begin{frame}{Probabilidad Condicional y Regla Total}
  \begin{columns}[T,onlytextwidth]
    \column{0.48\textwidth}
      \begin{block}{Probabilidad Condicional}
        ¿Cómo cambia la probabilidad de $B$ sabiendo que ocurrió $A$?
        \[
          P(B \mid A) = \frac{P(A \cap B)}{P(A)}
        \]
      \end{block}
      \small Si OptiBot sabe que el piso está mojado ($A$), la probabilidad de derrape ($B$) aumenta.
    \column{0.48\textwidth}
      \begin{alertblock}{Ley de Probabilidad Total}
        Si el espacio se divide en eventos disjuntos $A_1, A_2, \dots, A_n$:
        \[
          P(B) = \sum_{i=1}^n P(B \mid A_i) \cdot P(A_i)
        \]
      \end{alertblock}
      \small ¡Esta suma ponderada es el corazón exacto de las ecuaciones de Bellman!
  \end{columns}
  \vspace{0.4cm}
  \keyidea{La probabilidad condicional formaliza cómo la información actual condiciona las probabilidades del futuro.}
\end{frame}

\begin{frame}{El Concepto Rey del RL: Valor Esperado}
  \begin{columns}[c,onlytextwidth]
    \column{0.50\textwidth}
      \begin{beamercolorbox}[rounded=true,sep=0.8em]{block body}
        \centering
        {\Large $\mathbb{E}[X] = \sum_{x} x \cdot P(X = x)$}
      \end{beamercolorbox}
      \vspace{0.3cm}
      \begin{itemize}
        \item Es el promedio ponderado a largo plazo.
        \item \textbf{Cálculo manual en OptiBot:}\\
        Ganancia neta por paquete según tiempo $T$:
        \begin{itemize}
          \item 1 día ($+\$10$): prob $0.30$
          \item 2 días ($+\$5$): prob $0.50$
          \item 3 días ($-\$2$): prob $0.20$
        \end{itemize}
      \end{itemize}
    \column{0.46\textwidth}
      \begin{exampleblock}{Cálculo manual con Lápiz y Papel}
        \begin{align*}
          \mathbb{E}[\text{Ganancia}] &= (10)(0.30) + (5)(0.50) + (-2)(0.20) \\
          &= 3.0 + 2.5 - 0.4 \\
          &= \mathbf{\$5.10 \text{ por paquete}}
        \end{align*}
      \end{exampleblock}
      \small Aunque ningún paquete dé exactamente \$5.10, en 10,000 entregas ganaremos en promedio \$5.10 por paquete.
  \end{columns}
\end{frame}

\begin{frame}[fragile]{Verificación en Python: Ley de Grandes Números}
  \begin{columns}[T,onlytextwidth]
    \column{0.62\textwidth}
\begin{lstlisting}[style=rlcode]
import numpy as np

# Valores y distribucion teorica
ganancias = np.array([10.0, 5.0, -2.0])
probabilidades = np.array([0.30, 0.50, 0.20])

# Valor esperado teorico
E_teorico = np.sum(ganancias * probabilidades)

# Simulacion de 10,000 entregas de OptiBot
np.random.seed(42)
muestras = np.random.choice(ganancias, size=10000, p=probabilidades)
E_empirico = np.mean(muestras)

print(f"Teorico:  ${E_teorico:.2f}")
print(f"Empirico: ${E_empirico:.2f}")
# Salida: Teorico: $5.10 | Empirico: $5.08
\end{lstlisting}
    \column{0.35\textwidth}
      \begin{block}{Lección del código}
        \begin{itemize}
          \item Cuando no conocemos la distribución teórica, podemos \textbf{estimar la esperanza promediando muestras}.
          \item Este principio simple es la base de los métodos \textbf{Monte Carlo} y del aprendizaje por interacción.
        \end{itemize}
      \end{block}
  \end{columns}
\end{frame}

% ============================================================
% MÓDULO 1: BANDIDOS MULTI-BRAZO
% ============================================================
\section{M1 · Bandidos}
\module{Módulo 1: Bandidos Multi-Brazo}{Aprender por evaluación en un escenario de decisión estático}

\begin{frame}{OptiBot y la Selección de Proveedores}
  \begin{columns}[c,onlytextwidth]
    \column{0.52\textwidth}
      \textbf{El Escenario en la Estación de Despacho:}
      \begin{itemize}
        \item OptiBot debe despachar paquetes eligiendo entre \textbf{4 proveedores} de mensajería: $\{A, B, C, D\}$.
        \item Cada proveedor tiene tiempos y costos variables (estocásticos).
        \item Cada envío genera una recompensa neta $R_t$.
        \item \textbf{Problema:} No conocemos cuál proveedor es el mejor de antemano.
      \end{itemize}
    \column{0.45\textwidth}
      \centering
      \begin{tikzpicture}[x=1.1cm,y=1.1cm]
        \node[draw=Ink,fill=Mist!30,rounded corners=4pt,inner sep=6pt,align=center] (robot) at (0,1.2)
          {\textbf{OptiBot}\\{\scriptsize Estación fija}};
        \node[draw=Teal,fill=Teal!12,rounded corners=3pt] (p1) at (2.4,2.2) {Prov. A ($q_*=3.2$)};
        \node[draw=Coral,fill=Coral!12,rounded corners=3pt] (p2) at (2.4,1.4) {Prov. B ($q_*=4.8$)};
        \node[draw=Gold,fill=Gold!15,rounded corners=3pt] (p3) at (2.4,0.6) {Prov. C ($q_*=2.1$)};
        \node[draw=Leaf,fill=Leaf!15,rounded corners=3pt] (p4) at (2.4,-0.2) {Prov. D ($q_*=4.1$)};
        \draw[-{Latex},thick,SoftText] (robot) -- (p1);
        \draw[-{Latex},thick,Coral,line width=1.5pt] (robot) -- (p2);
        \draw[-{Latex},thick,SoftText] (robot) -- (p3);
        \draw[-{Latex},thick,SoftText] (robot) -- (p4);
      \end{tikzpicture}
  \end{columns}
  \vspace{0.3cm}
  \keyidea{En Bandidos sólo hay \textbf{acciones} y \textbf{recompensas}. No hay estados: el entorno no cambia de lugar.}
\end{frame}

\begin{frame}{Valor de una Acción y el Dilema Central}
  \begin{columns}[T,onlytextwidth]
    \column{0.48\textwidth}
      \begin{block}{Valor Verdadero $q_*(a)$}
        \[
          q_*(a) \doteq \mathbb{E}[R_t \mid A_t = a]
        \]
        Es la ganancia esperada condicionada a elegir la acción $a$.
      \end{block}
      \vspace{0.2cm}
      Nuestra estimación en el paso $t$ se denota $Q_t(a)$. Queremos que $Q_t(a) \to q_*(a)$.
    \column{0.48\textwidth}
      \begin{alertblock}{Exploración vs. Explotación}
        \begin{itemize}
          \item \textbf{Explotar:} Elegir la acción con mayor $Q_t(a)$ actual (maximizar ganancia inmediata).
          \item \textbf{Explorar:} Probar acciones no óptimas para aprender su verdadero $q_*(a)$ (ganancia futura).
        \end{itemize}
      \end{alertblock}
  \end{columns}
  \vspace{0.4cm}
  \keyidea{El agente nunca puede hacer ambas cosas al mismo tiempo con una sola acción: debe alternar con criterio.}
\end{frame}

\begin{frame}{La Regla de Actualización Incremental}
  \begin{columns}[c,onlytextwidth]
    \column{0.52\textwidth}
      ¿Cómo actualizamos el promedio de una acción sin guardar todo su historial en memoria?
      \[
        Q_{n+1} = Q_n + \alpha \Big( \underbrace{R_n - Q_n}_{\text{Error de Predicción}} \Big)
      \]
      \begin{itemize}
        \item Si $\alpha = \frac{1}{n}$, es el promedio muestral exacto.
        \item Si $\alpha = \text{cte} \in (0, 1]$, pondera más los resultados recientes (entornos no estacionarios).
      \end{itemize}
    \column{0.45\textwidth}
      \begin{beamercolorbox}[rounded=true,sep=0.7em]{block body}
        \centering\bfseries
        Nuevo $\leftarrow$ Viejo $+$ Paso $\times$ [Objetivo $-$ Viejo]
      \end{beamercolorbox}
      \vspace{0.3cm}
      \small \textbf{¡El patrón universal de RL!}\\
      Esta estructura reaparece en:
      \begin{itemize}
        \item Diferencia Temporal TD(0)
        \item SARSA y Q-Learning
        \item Gradiente descendente en Redes Neuronales
      \end{itemize}
  \end{columns}
\end{frame}

\begin{frame}[fragile]{Estrategia $\varepsilon$-Greedy en Python}
  \begin{columns}[T,onlytextwidth]
    \column{0.62\textwidth}
\begin{lstlisting}[style=rlcode]
import numpy as np

k = 4  # 4 proveedores
q_true = np.array([3.2, 4.8, 2.1, 4.1])  # desconocidos
Q = np.zeros(k)  # estimaciones iniciales
N = np.zeros(k)  # conteo de elecciones
epsilon = 0.10

for t in range(1, 1001):
    # Exploracion vs Explotacion
    if np.random.rand() < epsilon:
        action = np.random.choice(k)     # Explorar
    else:
        action = np.argmax(Q)            # Explotar (greedy)
    
    # Simular recompensa con ruido normal
    reward = np.random.normal(loc=q_true[action], scale=1.0)
    
    # Actualizacion incremental
    N[action] += 1
    Q[action] += (1.0 / N[action]) * (reward - Q[action])

print("Q estimado:", np.round(Q, 2))
# Salida convergente: [3.21, 4.79, 2.12, 4.09]
\end{lstlisting}
    \column{0.35\textwidth}
      \begin{exampleblock}{La gran limitación de M1}
        OptiBot aprende a elegir al proveedor $B$ óptimamente.\\
        \vspace{0.2cm}
        \textbf{Pero...} sus decisiones no cambian el mundo. El almacén siempre está en el mismo estado.\\
        \vspace{0.2cm}
        ¿Qué pasa cuando las acciones cambian el estado y el entorno evoluciona?
      \end{exampleblock}
  \end{columns}
\end{frame}

% ============================================================
% MÓDULO 2: CADENAS DE MARKOV
% ============================================================
\section{M2 · Cadenas de Markov}
\module{Módulo 2: Cadenas de Markov}{Modelando la evolución estocástica del entorno cuando no podemos intervenir}

\begin{frame}{OptiBot y la Cinta Transportadora}
  \begin{columns}[c,onlytextwidth]
    \column{0.50\textwidth}
      \textbf{El Escenario: Inspección de Rodillos}
      \begin{itemize}
        \item OptiBot vigila la cinta central de paquetes.
        \item Los rodamientos pasan por 3 estados:
        \begin{enumerate}
          \item \state{Nuevo ($N$)}
          \item \state{Desgastado ($D$)}
          \item \state{Averiado ($A$)}
        \end{enumerate}
        \item Al final de cada turno se inspeccionan.
        \item \textbf{Pregunta:} Si hoy está nuevo, ¿cuál es la probabilidad de que esté averiado en 2 turnos?
      \end{itemize}
    \column{0.48\textwidth}
      \centering
      \begin{tikzpicture}[scale=0.85,transform shape,
        st/.style={circle,draw=Ink,fill=Teal!12,thick,minimum size=1.1cm,font=\bfseries},
        edge/.style={-{Latex[length=2.5mm]},thick,SoftText}]
        \node[st] (N) at (0,2) {$N$};
        \node[st,fill=Gold!25] (D) at (2.5,2) {$D$};
        \node[st,fill=Coral!20] (A) at (1.25,0) {$A$};

        \draw[edge] (N) to[loop left] node[left] {$0.7$} (N);
        \draw[edge,Teal] (N) to[bend left=15] node[above] {$0.2$} (D);
        \draw[edge] (N) to[bend right=15] node[below left] {$0.1$} (A);

        \draw[edge] (D) to[loop right] node[right] {$0.6$} (D);
        \draw[edge,Coral] (D) to[bend left=15] node[below right] {$0.3$} (A);
        \draw[edge] (D) to[bend left=15] node[below] {$0.1$} (N);

        \draw[edge] (A) to[loop below] node[below] {$0.9$} (A);
        \draw[edge] (A) to[bend right=30] node[left] {$0.1$} (N);
      \end{tikzpicture}
  \end{columns}
  \vspace{0.3cm}
  \keyidea{La \textbf{Propiedad de Markov}: Dado el estado actual, el futuro es condicionalmente independiente del pasado.}
\end{frame}

\begin{frame}{La Matriz de Transición de Estados}
  \begin{columns}[c,onlytextwidth]
    \column{0.50\textwidth}
      \begin{block}{Definición Formal}
        Para cualquier secuencia de estados:
        \[
          P(S_{t+1} \mid S_t, S_{t-1}, \dots, S_0) = P(S_{t+1} \mid S_t)
        \]
      \end{block}
      \vspace{0.2cm}
      Agrupamos todas las probabilidades de transición en la matriz estocástica $P$:
      \[
        P_{ij} = P(S_{t+1} = j \mid S_t = i)
      \]
      \textbf{Regla de oro:} Cada fila debe sumar exactamente 1.
    \column{0.46\textwidth}
      \[
        P = \begin{pmatrix}
          0.7 & 0.2 & 0.1 \\
          0.1 & 0.6 & 0.3 \\
          0.1 & 0.0 & 0.9
        \end{pmatrix}
      \]
      \begin{itemize}
        \item Fila 1 ($N$): $0.7 + 0.2 + 0.1 = 1.0$
        \item Fila 2 ($D$): $0.1 + 0.6 + 0.3 = 1.0$
        \item Fila 3 ($A$): $0.1 + 0.0 + 0.9 = 1.0$
      \end{itemize}
      \vspace{0.2cm}
      \small ¿Probabilidad en 2 pasos? Simplemente calculamos $P^2 = P \times P$.
  \end{columns}
\end{frame}

\begin{frame}{Cálculo Manual: Distribución a 2 Pasos y Estado Estacionario}
  \begin{columns}[T,onlytextwidth]
    \column{0.48\textwidth}
      \begin{exampleblock}{Cálculo a 2 turnos desde $N$}
        Partiendo de $\mu_0 = [1, 0, 0]$:
        \begin{align*}
          \mu_1 &= \mu_0 P = [0.7, 0.2, 0.1] \\
          \mu_2 &= \mu_1 P \\
          P(A \text{ en } t=2) &= (0.7)(0.1) + (0.2)(0.3) + (0.1)(0.9) \\
          &= 0.07 + 0.06 + 0.09 = \mathbf{0.22}
        \end{align*}
      \end{exampleblock}
      \small Hay un $22\%$ de probabilidad de fallo al segundo turno.
    \column{0.48\textwidth}
      \begin{block}{Distribución Estacionaria $\pi^*$}
        A largo plazo, el sistema alcanza el equilibrio:
        \[
          \pi^* P = \pi^* \quad \text{con} \quad \sum_i \pi^*_i = 1
        \]
      \end{block}
      \vspace{0.2cm}
      \small Nos dice qué porcentaje del tiempo la cinta pasará averiada, sin importar dónde comenzó.
  \end{columns}
  \vspace{0.35cm}
  \keyidea{El estado estacionario describe la distribución de probabilidad invariante bajo la dinámica del sistema.}
\end{frame}

\begin{frame}[fragile]{Verificación Numérica en Python}
  \begin{columns}[T,onlytextwidth]
    \column{0.62\textwidth}
\begin{lstlisting}[style=rlcode]
import numpy as np

P = np.array([
    [0.7, 0.2, 0.1],
    [0.1, 0.6, 0.3],
    [0.1, 0.0, 0.9]
])

# Proyeccion a 2 turnos desde Nuevo
mu_0 = np.array([1.0, 0.0, 0.0])
mu_2 = mu_0 @ np.linalg.matrix_power(P, 2)
print("Estado en t=2:", np.round(mu_2, 3))
# [0.52 0.26 0.22] -> 22% de averia confirmada

# Estado estacionario: autovector con autovalor 1
evals, evecs = np.linalg.eig(P.T)
stationary = evecs[:, np.isclose(evals, 1.0)].real
stationary = stationary / np.sum(stationary)
print("Distribucion estacionaria pi*:", np.round(stationary.T[0], 3))
# [0.25 0.125 0.625]
\end{lstlisting}
    \column{0.35\textwidth}
      \begin{alertblock}{La limitación de M2}
        Ahora sabemos cómo evoluciona el mundo en el tiempo y podemos predecir el futuro.\\
        \vspace{0.2cm}
        \textbf{Pero somos espectadores pasivos}: OptiBot sólo mira la cinta romperse. No puede intervenir con acciones.\\
        \vspace{0.2cm}
        ¿Y si unimos las \textbf{acciones} de M1 con los \textbf{estados} de M2? $\to$ \textbf{¡MDP!}
      \end{alertblock}
  \end{columns}
\end{frame}

% ============================================================
% MÓDULO 3: PROCESOS DE DECISIÓN DE MARKOV Y BELLMAN
% ============================================================
\section{M3 · MDP + Bellman}
\module{Módulo 3: MDPs y Bellman}{La convergencia de estados, acciones y recompensas: el núcleo formal del RL}

\begin{frame}{El Gran Salto Conceptual: Uniendo M1 y M2}
  \centering
  \begin{tikzpicture}[node distance=1.2cm,
    card/.style={rounded corners=5pt,minimum width=3.8cm,minimum height=1.8cm,align=center,inner sep=8pt},
    arr/.style={-{Latex[length=3mm]},line width=1.4pt}]
    \node[card,draw=Coral,fill=Coral!10] (m1) {\textcolor{Coral}{\bfseries Módulo 1: Bandidos}\\[0.2em]Hay \textbf{acciones} y recompensas,\\pero \textbf{no hay estados}.};
    \node[card,draw=Gold!80!Ink,fill=Gold!12,right=1.5cm of m1] (m2) {\textcolor{Gold!70!Ink}{\bfseries Módulo 2: C. Markov}\\[0.2em]Hay \textbf{estados} y dinámica,\\pero \textbf{no hay acciones}.};
    \node[card,draw=Teal,fill=Teal!12,below=1.2cm of $(m1)!0.5!(m2)$,minimum width=7cm] (mdp)
      {\textcolor{TealDark}{\bfseries Módulo 3: Proceso de Decisión de Markov (MDP)}\\[0.2em]
       \textbf{Estados + Acciones + Recompensas + Horizonte temporal $\gamma$}};
    \draw[arr,Coral] (m1.south) -- (mdp.north west);
    \draw[arr,Gold!80!Ink] (m2.south) -- (mdp.north east);
  \end{tikzpicture}
  \vspace{0.35cm}
  \keyidea{En un MDP, las acciones del agente no sólo obtienen recompensas inmediatas, sino que \textbf{determinan el siguiente estado del mundo}.}
\end{frame}

\begin{frame}{La Tupla del MDP y el Retorno Descontado}
  \begin{columns}[T,onlytextwidth]
    \column{0.48\textwidth}
      \begin{block}{Definición Formal: $\langle \mathcal{S}, \mathcal{A}, \mathcal{P}, \mathcal{R}, \gamma \rangle$}
        \begin{itemize}
          \item $\mathcal{S}$: Conjunto de estados posibles.
          \item $\mathcal{A}$: Conjunto de acciones disponibles.
          \item $\mathcal{P}_{ss'}^a = P(S_{t+1}=s' \mid S_t=s, A_t=a)$.
          \item $\mathcal{R}(s,a) = \mathbb{E}[R_{t+1} \mid S_t=s, A_t=a]$.
          \item $\gamma \in [0, 1)$: Factor de descuento.
        \end{itemize}
      \end{block}
    \column{0.48\textwidth}
      \begin{alertblock}{El Retorno Acumulado $G_t$}
        \[
          G_t = R_{t+1} + \gamma R_{t+2} + \gamma^2 R_{t+3} + \cdots = \sum_{k=0}^{\infty} \gamma^k R_{t+k+1}
        \]
      \end{alertblock}
      \textbf{¿Por qué descontar con $\gamma$?}
      \begin{itemize}
        \item Evita infinitos matemáticos en tareas continuas.
        \item Modela la incertidumbre del futuro distante.
        \item Si $\gamma \to 0$, agente miope; si $\gamma \to 1$, visionario.
      \end{itemize}
  \end{columns}
\end{frame}

\begin{frame}{Funciones de Valor: $V^\pi(s)$ y $Q^\pi(s,a)$}
  \begin{columns}[T,onlytextwidth]
    \column{0.48\textwidth}
      \begin{block}{Función Valor de Estado $V^\pi(s)$}
        ¿Qué tan bueno es estar en el estado $s$ si sigo la política $\pi$?
        \[
          V^\pi(s) \doteq \mathbb{E}_\pi [G_t \mid S_t = s]
        \]
      \end{block}
      \vspace{0.2cm}
      \centering
      \begin{tikzpicture}[scale=0.85]
        \node[circle,draw=TealDark,fill=Teal!20,thick,minimum size=0.8cm] (s) at (0,1) {$s$};
        \node[circle,draw=Coral,fill=Coral!20,thick,minimum size=0.4cm] (a1) at (-0.8,0) {};
        \node[circle,draw=Coral,fill=Coral!20,thick,minimum size=0.4cm] (a2) at (0.8,0) {};
        \draw[-{Latex},thick] (s) -- node[left,font=\tiny] {$\pi(a_1|s)$} (a1);
        \draw[-{Latex},thick] (s) -- node[right,font=\tiny] {$\pi(a_2|s)$} (a2);
      \end{tikzpicture}
    \column{0.48\textwidth}
      \begin{alertblock}{Función Valor de Acción $Q^\pi(s,a)$}
        ¿Qué tan bueno es tomar $a$ en $s$, y a partir de ahí seguir $\pi$?
        \[
          Q^\pi(s,a) \doteq \mathbb{E}_\pi [G_t \mid S_t = s, A_t = a]
        \]
      \end{alertblock}
      \vspace{0.2cm}
      \centering
      \begin{tikzpicture}[scale=0.85]
        \node[circle,draw=Coral,fill=Coral!20,thick,minimum size=0.5cm] (sa) at (0,1) {$a$};
        \node[circle,draw=TealDark,fill=Teal!20,thick,minimum size=0.7cm] (s1) at (-0.8,0) {$s'_1$};
        \node[circle,draw=TealDark,fill=Teal!20,thick,minimum size=0.7cm] (s2) at (0.8,0) {$s'_2$};
        \draw[-{Latex},thick] (sa) -- node[left,font=\tiny] {$p$} (s1);
        \draw[-{Latex},thick] (sa) -- node[right,font=\tiny] {$p$} (s2);
      \end{tikzpicture}
  \end{columns}
\end{frame}

\begin{frame}{La Ecuación de Expectativa de Bellman}
  \begin{columns}[c,onlytextwidth]
    \column{0.55\textwidth}
      \begin{beamercolorbox}[rounded=true,sep=0.9em]{block body}
        \[
          V^\pi(s) = \sum_{a \in \mathcal{A}} \pi(a \mid s) \left[ \mathcal{R}(s,a) + \gamma \sum_{s' \in \mathcal{S}} \mathcal{P}_{ss'}^a V^\pi(s') \right]
        \]
      \end{beamercolorbox}
      \vspace{0.3cm}
      \textbf{La intuición fundamental de Richard Bellman:}
      \[
        \text{Valor de HOY} = \mathbb{E}[\text{Recompensa Inmediata}] + \gamma \times \mathbb{E}[\text{Valor de MAÑANA}]
      \]
    \column{0.42\textwidth}
      \begin{block}{Propiedad crucial}
        Bellman convierte un problema infinito y recurrente en un \textbf{sistema lineal de ecuaciones} con $|\mathcal{S}|$ incógnitas y $|\mathcal{S}|$ ecuaciones:
        \[
          \mathbf{V}^\pi = \mathbf{R}^\pi + \gamma \mathbf{P}^\pi \mathbf{V}^\pi
        \]
        cuya solución analítica es:
        \[
          \mathbf{V}^\pi = (\mathbf{I} - \gamma \mathbf{P}^\pi)^{-1} \mathbf{R}^\pi
        \]
      \end{block}
  \end{columns}
\end{frame}

\begin{frame}[fragile]{Cálculo Numérico en Python: Almacén OptiBot}
  \begin{columns}[T,onlytextwidth]
    \column{0.62\textwidth}
\begin{lstlisting}[style=rlcode]
import numpy as np

# Estados: 0: Muelle, 1: Pasillo, 2: Bateria
gamma = 0.90

# Dinamica bajo politica fija pi (ir a cargar)
P_pi = np.array([
    [0.1, 0.8, 0.1],  # Desde Muelle
    [0.0, 0.2, 0.8],  # Desde Pasillo
    [0.0, 0.0, 1.0]   # Bateria (estado terminal/carga)
])

# Recompensas esperadas inmediatas
R_pi = np.array([+2.0, -1.0, +10.0])

# Resolucion exacta de la Ecuacion de Bellman:
# V = (I - gamma * P_pi)^(-1) * R_pi
I = np.eye(3)
V_pi = np.linalg.inv(I - gamma * P_pi) @ R_pi

print("Valores V(s):", np.round(V_pi, 2))
# Salida: [77.58, 86.82, 100.0]
\end{lstlisting}
    \column{0.35\textwidth}
      \begin{alertblock}{La limitación de M3}
        Bellman nos permite \textbf{evaluar} el valor de los estados dada una política $\pi$ fija.\\
        \vspace{0.2cm}
        Pero OptiBot no quiere evaluar una política cualquiera: quiere \textbf{encontrar la política óptima $\pi^*$}.\\
        \vspace{0.2cm}
        ¿Cómo buscamos el comportamiento perfecto? $\to$ \textbf{Programación Dinámica (M4)}.
      \end{alertblock}
  \end{columns}
\end{frame}

% ============================================================
% MÓDULO 4: PROGRAMACIÓN DINÁMICA Y GPI
% ============================================================
\section{M4 · Prog. Dinámica}
\module{Módulo 4: Programación Dinámica y GPI}{Cálculo de la política óptima cuando conocemos las leyes completas del mundo}

\begin{frame}{Ecuaciones de Bellman de Optimalidad}
  \begin{columns}[c,onlytextwidth]
    \column{0.50\textwidth}
      \begin{block}{Definición de Política Óptima $\pi^*$}
        Una política $\pi^*$ es óptima si y sólo si $V^{\pi^*}(s) \ge V^\pi(s)$ para todo estado $s$ y toda política $\pi$.
      \end{block}
      \vspace{0.2cm}
      Las funciones de valor óptimas satisfacen el principio de optimalidad:
      \[
        V^*(s) = \max_{a \in \mathcal{A}} Q^*(s,a)
      \]
    \column{0.48\textwidth}
      \begin{beamercolorbox}[rounded=true,sep=0.8em]{block body alerted}
        \centering
        {\large $V^*(s) = \displaystyle\max_{a} \left[ \mathcal{R}(s,a) + \gamma \sum_{s'} \mathcal{P}_{ss'}^a V^*(s') \right]$}
      \end{beamercolorbox}
      \vspace{0.25cm}
      \small \textbf{¡Ojo con la no linealidad!}\\
      Debido al operador $\max$, este ya \textbf{no} es un sistema de ecuaciones lineales. No se resuelve con inversión de matrices directa. Requiere métodos iterativos.
  \end{columns}
\end{frame}

\begin{frame}{Policy Iteration vs. Value Iteration}
  \begin{columns}[T,onlytextwidth]
    \column{0.48\textwidth}
      \begin{block}{Policy Iteration (Howard, 1960)}
        Alterna dos fases explícitas hasta converger:
        \begin{enumerate}
          \item \textbf{Evaluación de Política:} Resolver $V^{\pi_k}$ hasta convergencia.
          \item \textbf{Mejora de Política:}
          \[
            \pi_{k+1}(s) \leftarrow \arg\max_a Q^{\pi_k}(s,a)
          \]
        \end{enumerate}
        \textbf{Ventaja:} Pocas iteraciones de mejora.\\
        \textbf{Desventaja:} Evaluar $V$ en cada paso es costoso.
      \end{block}
    \column{0.48\textwidth}
      \begin{alertblock}{Value Iteration (Bellman, 1957)}
        Combina evaluación y mejora en un solo paso:
        \[
          V_{k+1}(s) \leftarrow \max_a \left[ \mathcal{R}(s,a) + \gamma \sum_{s'} \mathcal{P}_{ss'}^a V_k(s') \right]
        \]
        \textbf{Garantía matemática:} El operador de Bellman es una \textbf{contracción} en norma infinito ($L_\infty$). Converge geométricamente a la única solución $V^*$.
      \end{alertblock}
  \end{columns}
\end{frame}

\begin{frame}{Interacción Política Generalizada (GPI)}
  \centering
  \begin{tikzpicture}[scale=0.95,transform shape]
    \draw[Teal,line width=2.5pt] (0,0) -- (7,3.5) node[right,font=\bfseries,text=Teal] {$v \approx v_\pi$};
    \draw[Coral,line width=2.5pt] (0,3.5) -- (7,0) node[right,font=\bfseries,text=Coral] {$\pi \approx \text{greedy}(v)$};

    % Trayectoria en zigzag de GPI
    \draw[-{Latex},line width=1.5pt,Ink] (0.8,0.4) -- (1.5,2.75) node[midway,left,font=\scriptsize] {evaluación};
    \draw[-{Latex},line width=1.5pt,Ink] (1.5,2.75) -- (2.6,1.3) node[midway,above right,font=\scriptsize] {mejora};
    \draw[-{Latex},line width=1.5pt,Ink] (2.6,1.3) -- (3.5,1.75);
    \draw[-{Latex},line width=1.5pt,Ink] (3.5,1.75) -- (4.8,1.1);

    \fill[Gold] (3.5,1.75) circle (0.15cm);
    \node[above=0.2cm of {(3.5,1.75)},font=\bfseries\small,text=Gold!80!Ink] {Punto óptimo $(\pi^*, v^*)$};
  \end{tikzpicture}
  \vspace{0.4cm}
  \keyidea{Casi todo el RL consiste en GPI: dos procesos compitiendo y colaborando que convergen al punto fijo óptimo.}
\end{frame}

\begin{frame}[fragile]{Implementación en Python: Value Iteration para OptiBot}
  \begin{columns}[T,onlytextwidth]
    \column{0.62\textwidth}
\begin{lstlisting}[style=rlcode]
import numpy as np

# Cuadricula simple 3x3 del almacen
n_states = 9
gamma = 0.95
V = np.zeros(n_states)
theta = 1e-4  # Umbral de convergencia

# Bucle de Value Iteration
while True:
    delta = 0
    for s in range(n_states):
        v = V[s]
        # Calcular valor para cada accion posible
        # Q(s, a) = R(s,a) + gamma * sum(P * V)
        q_values = []
        for a in range(4): # 0:Arriba, 1:Abajo, 2:Izq, 3:Der
            r = -1.0 if s != 8 else 0.0 # Recompensa
            s_next = s # Logica de transicion simple
            q_values.append(r + gamma * V[s_next])
        V[s] = max(q_values)
        delta = max(delta, abs(v - V[s]))
    if delta < theta:
        break
\end{lstlisting}
    \column{0.35\textwidth}
      \begin{alertblock}{La limitación de M4}
        DP es elegante y exacto.\\
        \vspace{0.2cm}
        \textbf{Pero exige conocer $\mathcal{P}$ y $\mathcal{R}$ de antemano.}\\
        \vspace{0.2cm}
        En un almacén real con obstáculos móviles y suelo irregular, \textbf{nadie} nos entrega la matriz de transición.\\
        \vspace{0.2cm}
        ¿Cómo aprende OptiBot sólo interactuando por ensayo y error? $\to$ \textbf{Monte Carlo (M5)}.
      \end{alertblock}
  \end{columns}
\end{frame}

% ============================================================
% MÓDULO 5: MÉTODOS MONTE CARLO
% ============================================================
\section{M5 · Monte Carlo}
\module{Módulo 5: Métodos Monte Carlo}{Aprender de la experiencia sin modelo: promediando historias completas}

\begin{frame}{El Cambio de Paradigma: Model-Free RL}
  \begin{columns}[c,onlytextwidth]
    \column{0.50\textwidth}
      \textbf{El nuevo escenario de OptiBot:}
      \begin{itemize}
        \item No tenemos planos ni fórmulas matemáticas de transición.
        \item El robot simplemente ejecuta un turno completo (\textbf{episodio}).
        \item Registra la historia observada:
        \[
          S_0, A_0, R_1, S_1, A_1, R_2, \dots, S_T
        \]
        \item Al sonar la campana de fin de turno, calcula cuánto ganó a partir de cada estado visitado.
      \end{itemize}
    \column{0.48\textwidth}
      \begin{beamercolorbox}[rounded=true,sep=0.8em]{block body}
        \centering
        {\large $V(s) \approx \frac{1}{N(s)} \sum_{i=1}^{N(s)} G_i(s)$}
      \end{beamercolorbox}
      \vspace{0.3cm}
      \small Recordando el \textbf{Módulo 0}: El promedio muestral de retornos empíricos converge exactamente a la esperanza matemática $\mathbb{E}[G_t \mid S_t=s]$ por la Ley de Grandes Números.
  \end{columns}
  \vspace{0.4cm}
  \keyidea{Monte Carlo no requiere conocer la física del entorno; sólo necesita \textbf{experiencia simulada o real}.}
\end{frame}

\begin{frame}{First-Visit MC y el Retorno hacia Atrás}
  \begin{columns}[T,onlytextwidth]
    \column{0.52\textwidth}
      \begin{exampleblock}{Cálculo Manual Paso a Paso}
        Trayectoria observada con $\gamma = 0.9$:
        \begin{center}
          $S_0 \xrightarrow{r_1=-1} S_1 \xrightarrow{r_2=-1} S_2 \xrightarrow{r_3=+10} \text{Meta}$
        \end{center}
        Calculamos los retornos acumulados \textbf{de atrás hacia adelante}:
        \begin{align*}
          G_2 &= r_3 = \mathbf{+10.0} \\
          G_1 &= r_2 + \gamma G_2 = -1 + 0.9(10.0) = \mathbf{+8.0} \\
          G_0 &= r_1 + \gamma G_1 = -1 + 0.9(8.0) = \mathbf{+6.2}
        \end{align*}
      \end{exampleblock}
    \column{0.45\textwidth}
      \begin{block}{First-Visit vs. Every-Visit}
        \begin{itemize}
          \item \textbf{First-Visit:} Se promedia $G_t$ sólo la primera vez que se visita el estado $s$ en el episodio. Estimador insesgado.
          \item \textbf{Every-Visit:} Se promedia cada vez que se toca $s$. Ligeramente sesgado pero menor error cuadrático inicial.
        \end{itemize}
      \end{block}
      \vspace{0.2cm}
      \small Para control, estimamos $Q(s,a)$ para poder elegir la acción codiciosa sin requerir modelo.
  \end{columns}
\end{frame}

\begin{frame}[fragile]{Control Monte Carlo $\varepsilon$-Greedy en Python}
  \begin{columns}[T,onlytextwidth]
    \column{0.62\textwidth}
\begin{lstlisting}[style=rlcode]
# Registro acumulado de retornos para cada par (s, a)
returns = {(s, a): [] for s in states for a in actions}
Q = np.zeros((n_states, n_actions))

for ep in range(n_episodes):
    trajectory = run_episode(policy, env)
    G = 0
    visited = set()
    
    # Recorrido de atras hacia adelante
    for s, a, r in reversed(trajectory):
        G = gamma * G + r
        if (s, a) not in visited:
            visited.add((s, a))
            returns[(s, a)].append(G)
            # Actualizacion de la tabla Q
            Q[s, a] = np.mean(returns[(s, a)])
            
    # Mejora de politica epsilon-greedy
    for s in states:
        best_a = np.argmax(Q[s])
        # Actualizar probabilidades de exploracion...
\end{lstlisting}
    \column{0.35\textwidth}
      \begin{alertblock}{La limitación de M5}
        Monte Carlo es insesgado y no necesita modelo.\\
        \vspace{0.2cm}
        \textbf{Pero tiene 2 problemas graves:}
        \begin{enumerate}
          \item \textbf{Alta varianza:} Un solo giro al final altera todos los retornos.
          \item \textbf{Obligatoriamente episódico:} Debe esperar a que termine el día completo para actualizar el primer paso.
        \end{enumerate}
        \vspace{0.2cm}
        ¿Podemos aprender \textbf{paso a paso} sin esperar al final? $\to$ \textbf{TD Learning (M6)}.
      \end{alertblock}
  \end{columns}
\end{frame}

% ============================================================
% MÓDULO 6: TEMPORAL-DIFFERENCE LEARNING (SARSA Y Q-LEARNING)
% ============================================================
\section{M6 · TD Learning}
\module{Módulo 6: TD Learning}{La combinación definitiva: muestreo sin modelo + estimación por bootstrapping}

\begin{frame}{La Idea Central de la Diferencia Temporal: Bootstrapping}
  \begin{columns}[c,onlytextwidth]
    \column{0.50\textwidth}
      \textbf{¿Qué es Bootstrapping?}
      Actualizar una predicción a partir de \textbf{otra predicción futura}, sin esperar al desenlace final del episodio.
      \vspace{0.25cm}
      \begin{itemize}
        \item \textbf{DP:} Hace bootstrapping, pero necesita el modelo $\mathcal{P}$ y $\mathcal{R}$.
        \item \textbf{MC:} No hace bootstrapping (espera al final), pero no necesita modelo.
        \item \textbf{TD:} \textbf{¡Lo mejor de ambos mundos!} Aprende de la experiencia pura y actualiza paso a paso.
      \end{itemize}
    \column{0.48\textwidth}
      \begin{beamercolorbox}[rounded=true,sep=0.8em]{block body}
        \centering
        {\large $V(S_t) \leftarrow V(S_t) + \alpha \big[ \underbrace{R_{t+1} + \gamma V(S_{t+1})}_{\text{Objetivo TD}} - V(S_t) \big]$}
      \end{beamercolorbox}
      \vspace{0.25cm}
      \small El \textbf{Error TD} $\delta_t = R_{t+1} + \gamma V(S_{t+1}) - V(S_t)$ mide la sorpresa entre lo que predije ayer y la nueva evidencia de hoy.
  \end{columns}
  \vspace{0.35cm}
  \keyidea{TD actualiza en cada paso individual ($t \to t+1$): puede aprender en tareas continuas sin fin de turno.}
\end{frame}

\begin{frame}{Control: SARSA vs. Q-Learning}
  \begin{columns}[T,onlytextwidth]
    \column{0.48\textwidth}
      \begin{block}{SARSA (On-Policy)}
        Aprende el valor de la política que realmente ejecuta (incluyendo los errores de exploración):
        \[
          Q(S,A) \leftarrow Q(S,A) + \alpha [ R + \gamma Q(S', A') - Q(S,A) ]
        \]
        \begin{itemize}
          \item La tupla: $(S_t, A_t, R_{t+1}, S_{t+1}, A_{t+1})$.
          \item $A'$ se elige con la misma política $\varepsilon$-greedy.
          \item Tiende a ser más \textbf{prudente y conservador}.
        \end{itemize}
      \end{block}
    \column{0.48\textwidth}
      \begin{alertblock}{Q-Learning (Off-Policy, Watkins 1989)}
        Aprende la política óptima $Q^*$ sin importar qué política de exploración use el robot:
        \[
          Q(S,A) \leftarrow Q(S,A) + \alpha [ R + \gamma \max_{a'} Q(S', a') - Q(S,A) ]
        \]
        \begin{itemize}
          \item El objetivo usa $\max_{a'}$, no la acción ejecutada.
          \item Aprende la ruta \textbf{óptima más corta}.
        \end{itemize}
      \end{alertblock}
  \end{columns}
\end{frame}

\begin{frame}{Cálculo Manual Paso a Paso: Actualización de Q-Learning}
  \begin{columns}[c,onlytextwidth]
    \column{0.50\textwidth}
      \textbf{Transición observada de OptiBot:}
      \begin{itemize}
        \item Estado actual: $S = \text{Pasillo A}$
        \item Acción: $A = \text{Avanzar}$
        \item Recompensa recibida: $R = -1.0$
        \item Siguiente estado: $S' = \text{Pasillo B}$
      \end{itemize}
      \textbf{Valores previos en la tabla Q:}
      \begin{itemize}
        \item $Q(S, A) = 4.0$
        \item En $S'$: $Q(S', \text{Izq}) = 6.0$, $Q(S', \text{Der}) = 8.0$
        \item Por tanto: $\max_{a'} Q(S', a') = 8.0$
        \item Parámetros: $\alpha = 0.20$, $\gamma = 0.90$
      \end{itemize}
    \column{0.48\textwidth}
      \begin{exampleblock}{Cálculo en la pizarra}
        \begin{align*}
          \text{Objetivo TD} &= R + \gamma \max_{a'} Q(S', a') \\
          &= -1.0 + 0.90(8.0) = \mathbf{6.20} \\
          \delta_t &= \text{Objetivo} - Q(S, A) \\
          &= 6.20 - 4.00 = \mathbf{+2.20} \\
          Q_{\text{nuevo}}(S, A) &= 4.00 + 0.20(+2.20) \\
          &= 4.00 + 0.44 = \mathbf{4.44}
        \end{align*}
      \end{exampleblock}
      \small ¡La estimación subió de 4.00 a 4.44 al descubrir que el Pasillo B tiene un camino muy prometedor!
  \end{columns}
\end{frame}

\begin{frame}[fragile]{Implementación de Q-Learning para el Almacén}
  \begin{columns}[T,onlytextwidth]
    \column{0.62\textwidth}
\begin{lstlisting}[style=rlcode]
for ep in range(num_episodes):
    s, _ = env.reset()
    done = False
    
    while not done:
        # Seleccion de accion epsilon-greedy
        if np.random.rand() < epsilon:
            a = env.action_space.sample()
        else:
            a = np.argmax(Q[s])
            
        s_next, r, terminated, truncated, _ = env.step(a)
        done = terminated or truncated
        
        # Actualizacion de Q-Learning
        target = r
        if not done:
            target += gamma * np.max(Q[s_next])
            
        td_error = target - Q[s, a]
        Q[s, a] += alpha * td_error
        s = s_next
\end{lstlisting}
    \column{0.35\textwidth}
      \begin{alertblock}{La limitación de M6}
        Q-Learning tabular es el estándar de oro para problemas pequeños.\\
        \vspace{0.2cm}
        \textbf{El abismo de la dimensionalidad:}
        Si el almacén tiene $1,000,000$ posiciones o sensores continuos (LiDAR, batería en \%, velocidades), la tabla requeriría gigabytes y nunca visitaría cada casilla.\\
        \vspace{0.2cm}
        ¿Cómo generalizar a estados nunca vistos? $\to$ \textbf{Deep Q-Networks (M7)}.
      \end{alertblock}
  \end{columns}
\end{frame}

% ============================================================
% MÓDULO 7: DEEP Q-NETWORKS (DQN)
% ============================================================
\section{M7 · Deep Q-Networks}
\module{Módulo 7: Deep Q-Networks}{Sustituyendo la tabla por redes neuronales: estabilidad, replay buffer y target network}

\begin{frame}{El Salto a la Aproximación de Funciones}
  \begin{columns}[c,onlytextwidth]
    \column{0.50\textwidth}
      \textbf{OptiBot con Sensores Continuos:}
      \begin{itemize}
        \item Estado continuo $s \in \mathbb{R}^4$:
        \begin{enumerate}
          \item Distancia frontal LiDAR ($0.1$ a $15.0$ m)
          \item Porcentaje de batería ($0.0\%$ a $100.0\%$)
          \item Velocidad lineal ($0.0$ a $2.5$ m/s)
          \item Ángulo respecto a la estación de carga
        \end{enumerate}
        \item Hay infinitos estados posibles: \textbf{una tabla Q es inviable}.
        \item \textbf{Solución:} Una red neuronal paramétrica $Q(s, a; \theta) \approx Q^*(s, a)$.
      \end{itemize}
    \column{0.48\textwidth}
      \centering
      \begin{tikzpicture}[scale=0.85,transform shape]
        % Capa entrada
        \foreach \i in {1,2,3,4} {
          \node[circle,draw=Teal,fill=Teal!20,thick,minimum size=0.55cm] (in\i) at (0, 2.5 - 0.9*\i) {\tiny $s_\i$};
        }
        % Capa oculta
        \foreach \j in {1,2,3,4,5} {
          \node[circle,draw=Ink,fill=Mist!40,thick,minimum size=0.55cm] (h\j) at (2.0, 3.0 - 0.8*\j) {};
        }
        % Capa salida
        \foreach \k/\act in {1/Avanzar, 2/Girar, 3/Cargar} {
          \node[circle,draw=Coral,fill=Coral!20,thick,minimum size=0.55cm] (out\k) at (4.0, 2.1 - 0.9*\k) {};
          \node[right=0.1cm of out\k,font=\tiny] {$Q(s,\text{\act})$};
        }
        \foreach \i in {1,2,3,4} {
          \foreach \j in {1,2,3,4,5} {
            \draw[-{Latex[length=1.5mm]},SoftText!40] (in\i) -- (h\j);
          }
        }
        \foreach \j in {1,2,3,4,5} {
          \foreach \k in {1,2,3} {
            \draw[-{Latex[length=1.5mm]},Coral!40] (h\j) -- (out\k);
          }
        }
      \end{tikzpicture}
  \end{columns}
  \vspace{0.25cm}
  \keyidea{La red neuronal permite \textbf{generalizar}: estados similares obtienen valores Q similares aunque nunca antes hayan sido visitados.}
\end{frame}

\begin{frame}{La Tríada Mortal y los Dos Pilares de DQN}
  \begin{columns}[T,onlytextwidth]
    \column{0.46\textwidth}
      \begin{alertblock}{La Tríada Mortal (Sutton \& Barto)}
        Causa divergencia si se combinan:
        \begin{enumerate}
          \item \textbf{Aproximación de funciones}
          \item \textbf{Bootstrapping} (objetivo TD)
          \item \textbf{Entrenamiento Off-policy}
        \end{enumerate}
      \end{alertblock}
      \small Para evitar que la red diverja, DeepMind introdujo dos mecanismos estabilizadores fundamentales.
    \column{0.50\textwidth}
      \begin{block}{1. Experience Replay (Buffer $\mathcal{D}$)}
        \begin{itemize}
          \item Guarda tuplas pasadas $(s, a, r, s', d)$.
          \item Muestrea mini-batches aleatorios uniformes.
          \item Rompe la correlación temporal entre pasos consecutivos.
        \end{itemize}
      \end{block}
      \begin{block}{2. Target Network (Red Objetivo $\theta^-$)}
        \begin{itemize}
          \item Congela los pesos del objetivo durante $C$ pasos.
          \item Evita perseguir un ``blanco móvil''.
        \end{itemize}
      \end{block}
  \end{columns}
\end{frame}

\begin{frame}{Función de Pérdida y Double DQN}
  \begin{columns}[c,onlytextwidth]
    \column{0.52\textwidth}
      \textbf{Función de Pérdida de DQN:}
      \[
        \mathcal{L}(\theta) = \mathbb{E}_{(s,a,r,s',d) \sim \mathcal{D}} \left[ \Big( y_i - Q(s,a;\theta) \Big)^2 \right]
      \]
      con el objetivo de Bellman:
      \[
        y_i = r + \gamma (1 - d) \max_{a'} Q(s', a'; \mathbf{\theta^-})
      \]
    \column{0.45\textwidth}
      \begin{exampleblock}{Mejora: Double DQN}
        El operador $\max$ induce un \textbf{sesgo de sobreestimación} sistemático. Double DQN desacopla la decisión de la valoración:
        \[
          y_i^{\text{DDQN}} = r + \gamma Q\Big( s', \arg\max_{a'} Q(s', a'; \mathbf{\theta}); \mathbf{\theta^-} \Big)
        \]
        \begin{itemize}
          \item La red online $\theta$ elige la acción.
          \item La red target $\theta^-$ evalúa su valor.
        \end{itemize}
      \end{exampleblock}
  \end{columns}
\end{frame}

\begin{frame}[fragile]{Estructura de DQN en PyTorch}
  \begin{columns}[T,onlytextwidth]
    \column{0.64\textwidth}
\begin{lstlisting}[style=rlcode]
import torch
import torch.nn as nn

class QNetwork(nn.Module):
    def __init__(self, state_dim, action_dim):
        super().__init__()
        self.net = nn.Sequential(
            nn.Linear(state_dim, 64),
            nn.ReLU(),
            nn.Linear(64, 64),
            nn.ReLU(),
            nn.Linear(64, action_dim)
        )
    def forward(self, x):
        return self.net(x)

# Calculo de perdida con Replay Buffer
states, actions, rewards, next_states, dones = replay_buffer.sample(32)
with torch.no_grad():
    # Objetivo con Target Network congelada
    max_next_q = target_net(next_states).max(1)[0].unsqueeze(1)
    target = rewards + gamma * (1 - dones) * max_next_q

current_q = online_net(states).gather(1, actions)
loss = nn.MSELoss()(current_q, target)
\end{lstlisting}
    \column{0.33\textwidth}
      \begin{block}{Entrenamiento Estable}
        \begin{enumerate}
          \item Muestrear batch.
          \item Calcular target fijo.
          \item Backpropagation en $\theta$.
          \item Actualizar periódicamente $\theta^- \leftarrow \theta$.
        \end{enumerate}
      \end{block}
  \end{columns}
\end{frame}

% ============================================================
% MÓDULO 8: PROYECTO FINAL Y FRONTERAS DEL RL
% ============================================================
\section{M8 · Proyecto y Fronteras}
\module{Módulo 8: Proyecto Integrador y Fronteras}{Diagnóstico experimental, diseño de recompensas y el camino hacia el RL moderno}

\begin{frame}{El Proyecto Final: OptiBot en el Almacén Dinámico}
  \begin{columns}[c,onlytextwidth]
    \column{0.50\textwidth}
      \textbf{El Desafío de Integración:}
      \begin{itemize}
        \item Ponemos a prueba a OptiBot en un entorno que cambia en tiempo real:
        \begin{itemize}
          \item Estantes móviles y pasillos bloqueados.
          \item Zonas de batería baja y estaciones de recarga.
          \item Pedidos urgentes con ventanas de entrega estrictas.
        \end{itemize}
        \item \textbf{Evaluación Comparativa:}
        \begin{enumerate}
          \item Política Aleatoria (Línea base)
          \item Q-Learning Tabular (Módulo 6)
          \item Deep Q-Network (Módulo 7)
        \end{enumerate}
      \end{itemize}
    \column{0.48\textwidth}
      \centering
      \begin{tikzpicture}[x=0.075cm,y=0.050cm]
        \draw[-{Latex},SoftText] (0,0) -- (55,0) node[right,font=\tiny] {Episodios};
        \draw[-{Latex},SoftText] (0,-10) -- (0,55) node[above,font=\tiny] {Recompensa media};
        % Aleatorio
        \draw[SoftText,dashed,line width=1pt] (0,-5) -- (50,-5) node[right,font=\tiny] {Aleatorio};
        % Q-Learning Tabular
        \draw[Gold!90!Ink,line width=1.5pt] (0,-8) .. controls (10,5) and (25,25) .. (50,30) node[right,font=\tiny] {Tabular};
        % DQN
        \draw[Teal,line width=2pt] (0,-8) .. controls (15,-5) and (25,35) .. (50,48) node[right,font=\tiny] {DQN};
      \end{tikzpicture}
  \end{columns}
  \vspace{0.3cm}
  \keyidea{El agente DQN demuestra capacidad de adaptación ante configuraciones no vistas durante el entrenamiento.}
\end{frame}

\begin{frame}{Diagnóstico y Patologías Comunes en RL}
  \begin{columns}[T,onlytextwidth]
    \column{0.48\textwidth}
      \begin{alertblock}{¿Por qué no aprende mi agente?}
        \begin{itemize}
          \item \textbf{Exploración insuficiente:} $\varepsilon$ decayó demasiado rápido; el agente nunca descubrió la meta.
          \item \textbf{Recompensas demasiado dispersas (sparse):} Si sólo hay $+100$ al final de 1,000 pasos, el gradiente no fluye.
          \item \textbf{Tasa de aprendizaje inestable:} Pérdida que oscila salvajemente o diverge a $\infty$.
        \end{itemize}
      \end{alertblock}
    \column{0.48\textwidth}
      \begin{block}{El Peligro del Reward Hacking}
        Si recompensamos a OptiBot por ``no colisionar'', el robot puede decidir \textbf{quedarse completamente quieto}.\\
        \vspace{0.2cm}
        \textbf{Lección de diseño:}
        El agente no maximiza lo que tú \emph{deseas} que haga; maximiza estricta y despiadadamente la función matemática que \emph{programaste}.
      \end{block}
  \end{columns}
\end{frame}

\begin{frame}{Fronteras: Del Control Discreto al RL Moderno}
  \begin{columns}[T,onlytextwidth]
    \column{0.32\textwidth}
      \begin{block}{Policy Gradients}
        \textbf{REINFORCE y PPO}\\
        \vspace{0.15cm}
        Optimizan directamente la política $\pi_\theta(a \mid s)$ sin tabla Q. Esencial para controlar brazos robóticos en espacios continuos.
      \end{block}
    \column{0.32\textwidth}
      \begin{alertblock}{Actor-Critic}
        \textbf{A2C, SAC, TD3}\\
        \vspace{0.15cm}
        Combinan un Actor (política que decide) con un Crítico (función de valor que evalúa el error TD). Máxima eficiencia muestral.
      \end{alertblock}
    \column{0.32\textwidth}
      \begin{exampleblock}{RLHF en LLMs}
        \textbf{Alineación de Modelos}\\
        \vspace{0.15cm}
        El aprendizaje por refuerzo con feedback humano es el motor que alinea ChatGPT, Claude y Gemini hacia respuestas útiles y seguras.
      \end{exampleblock}
  \end{columns}
  \vspace{0.4cm}
  \keyidea{Todos estos avances avanzados se fundamentan en los mismos tres pilares que dominaste: \textbf{Probabilidad, Bellman y Bootstrapping}.}
\end{frame}


% ============================================================
% CIERRE
% ============================================================
\begin{frame}[plain]
  \begin{tikzpicture}[remember picture,overlay]
    \fill[Paper] (current page.south west) rectangle (current page.north east);
    \fill[Teal] (current page.south west) rectangle ([yshift=0.13\paperheight]current page.south east);
    \draw[Coral,line width=3pt] ([xshift=-2.2cm,yshift=1.7cm]current page.center)
      -- ([xshift=2.2cm,yshift=1.7cm]current page.center);
    \node[text=Ink,font=\fontsize{26}{30}\selectfont\bfseries] at ([yshift=0.65cm]current page.center) {¡Felicitaciones! Has dominado el RL};
    \node[text=SoftText,font=\large] at ([yshift=-0.45cm]current page.center) {De la probabilidad elemental a las redes neuronales profundas.};
    \node[text=white,font=\small] at ([yshift=0.08\paperheight]current page.south) {OptiBot Logistics · Código Abierto y Reproducible};
  \end{tikzpicture}
\end{frame}

\end{document}
